% ------------------------------------------------------------------------
% ------------------------------------------------------------------------
% ------------------------------------------------------------------------
%                                Resumen
% ------------------------------------------------------------------------
% ------------------------------------------------------------------------
% ------------------------------------------------------------------------
\chapter*{RESUMEN}

\footnotesize{
  \begin{description}
    \item[TÍTULO:] \MakeUppercase{\titulo}
          \astfootnote{Trabajo de grado}
    \item[AUTOR:]{\autora}, {\autorb}
          \asttfootnote{\facultad. \escuela. Director: \director. Codirectores:
          \codirectora}
    \item[PALABRAS CLAVE:] RETINOPATÍA DIABÉTICA, REDES NEURONALES CONVOLUCIONALES,
    APRENDIZAJE PROFUNDO, IMÁGENES DEL FONDO DEL OJO, PROCESAMIENTO DIGITAL DE IMÁGENES,
    APRENDIZAJE POR TRANSFERENCIA.
    \item[DESCRIPCIÓN:] \hfill \\
          La retinopatía diabética (RD) es una complicación microvascular causada por
          daños en los vasos sanguíneos de la retina debido a niveles elevados de azúcar en sangre,
          común en personas con diabetes y es considerada una de las principales causas de
          ceguera en el mundo. Esta enfermedad suele diagnosticarse mediante exámenes oculares en
          los que se evalúa el fondo del ojo. Sin embargo, el diagnóstico manual es desafiante
          porque es costoso y consume tiempo, lo que limita su acceso en áreas con pocos
          recursos médicos. Esto hace necesario buscar métodos automatizados que permitan
          una detección más rápida y precisa de la enfermedad. Este trabajo desarrolla modelos 
          basados en redes neuronales convolucionales (CNN) para detectar la enfermedad y clasificar
          sus niveles de severidad mediante clasificadores binarios (presencia/ausencia) y multiclase
          (cinco niveles). Se utilizó aprendizaje por transferencia con las arquitecturas MobileNetV2, 
          ResNet50V2, InceptionV3 y VGG19, entrenadas con 35.121 imágenes del dataset público de Kaggle. 
          El preprocesamiento incluyó normalización de color, enmascaramiento periférico, balanceo de 
          clases mediante submuestreo y aumento de datos, incorporando bloques Multi-Head Attention, 
          capas GaussianNoise y técnicas de regularización (Dropout, L2, early stopping). El modelo 
          fue evaluado mediante sensibilidad, precisión, F1-score y AUC-ROC, siendo MobileNetV2 el 
          modelo de mejor desempeño, alcanzando una sensibilidad de 96,5\% en clasificación binaria 
          con intervalos de confianza (IC95\%: AUC-ROC 0.9636 - 0.99). Para clasificación multiclase,
          se alcanzó una sensibilidad de 65\%, identificando el canal verde como el más efectivo debido
          a su contraste para lesiones hemorrágicas. La validación externa en datasets MESSIDOR-2, 
          IDRiD y FOSCAL confirmó la capacidad de generalización del modelo, demostrando su viabilidad 
          como herramienta de apoyo al tamizaje clínico de retinopatía diabética.      
  \end{description}}\normalsize
% ------------------------------------------------------------------------ 